\documentclass[../../../../main.tex]{subfiles}
\begin{document}

\begin{flushright}
\footnotesize\textit{【自動文字起こし・要確認】}
\end{flushright}

[解] (1) から $|z_k|^p = 1 \ \therefore |z_k| = 1 \cdots \text{\textcircled{1}}$。又 $z_k \neq 1 \cdots \text{\textcircled{2}}$
\textcircled{1}、\textcircled{2}から $z_k = \cos\theta + i\sin\theta\ (0 < \theta < 2\pi)$ とおくと、$z_k = e(i\theta_k)\ (0 < \theta_k < 2\pi)$ とおける。(イ)、(ロ)から
\begin{eqnarray*}
\begin{cases} e(p\theta_k) = 1 \\ e(\frac{2\pi}{p}\theta_k) = 1 \end{cases} \cdots \text{\textcircled{3}}
\end{eqnarray*}
したがって、\textcircled{3}を満たす $\{\theta_k\}$ の個数が $a_n$ である。\textcircled{3}の第1式から、適当な自然数 $m_k$ を用いて
\[ \theta_k = \frac{2\pi}{p}m_k\ (m_k = 1, 2, \cdots, p-1) \cdots \text{\textcircled{4}} \]
とおける。したがって、$\{\theta_k\}$ の個数は $\{m_k\}$ の数に他ならない。\textcircled{3}の第2式に\textcircled{4}を代入して整理して、
\[ \frac{2\pi}{p} m_k \equiv 0 \pmod p \cdots \text{\textcircled{5}} \]
又、題意の $n$ を $n \in \mathbb{N}$ に拡張してかんがえる。
\[ a_1 = 0,\ a_2 = p-1,\ a_3 = (p-1)(p-2) \cdots (1) \]
(2) $S_n = \sum_{k=1}^n m_k$ とおく。$m_k = 1, 2, \cdots, p-1$ から
\[ \begin{cases} S_n \equiv 0 \pmod p \text{ の時、} S_{n+1} \not\equiv 0 \\ S_n \not\equiv 0 \pmod p \text{ の時、対応する } m_{n+1} \text{ がただ } 1 \text{ つあって、} S_{n+1} \equiv 0 \text{ となる。} \end{cases} \]
$S_n$ の値のとり方は、$(p-1)^n$ 通りだからしたがって、$S_n \equiv 0$ となる場合の数は
\[ a_{n+1} = (p-1)^n - a_n \]
(3) $b_n = \frac{a_n}{(p-1)^n}$ とおく。(2)から
\begin{eqnarray*}
b_{n+1} &=& -\frac{1}{p-1} b_n + \frac{1}{p-1} \\
b_{n+1} - \frac{1}{p} &=& -\frac{1}{p-1} \left(b_n - \frac{1}{p}\right)
\end{eqnarray*}
$b_1=0$ だから、くり返し用いて、
\begin{eqnarray*}
b_n &=& \left(-\frac{1}{p-1}\right)^{n-1} \left(-\frac{1}{p}\right) + \frac{1}{p} \\
a_n &=& \frac{1}{p} \left[ (p-1)^n + (-1)^n (p-1) \right]
\end{eqnarray*}

[解2] (法$p$とする)
(2) $a_{n+2}$ についてかんがえる。
$1^\circ\ m_{n+1} + m_{n+2} \equiv 0$ の時
このような $(m_{n+1}, m_{n+2})$ のくみあわせは $1 \le m_i \le p-1$ から、$m_{n+1}$ を定めればただ $1$ 通りに定まり、$p-1$ 通り。又、$S_{n+2} \equiv 0 \iff S_n \equiv 0$ だから、$m_1 \sim m_n$ の定め方は $a_n$ 通り。あわせて
\[ (p-1) a_n \text{ 通り} \]
$2^\circ\ m_{n+1} + m_{n+2} \equiv r\ (r=1, 2, \cdots, p-1)$ の時
$m_{n+1}$ のえらび方は $1, 2, \cdots, p-1$ から $r$ をのぞいた $p-2$ 通り。
$m_{n+2}$ は、
\begin{eqnarray*}
m_{n+1} < r \text{ ならば } m_{n+2} &=& r - m_{n+1} \\
m_{n+1} > r \text{ ならば } m_{n+2} &=& p + r - m_{n+1}
\end{eqnarray*}
と各々 $1$ 通りに定まる。又、$S_{n+2} \equiv 0 \iff S_n + r \equiv 0$ だから、$m_1 \sim m_n$ の定め方は $a_{n+1}$ 通り。あわせて
\[ (p-2) a_{n+1} \]
以上で全ての場合がつくされ、$1^\circ, 2^\circ$ は排反だから
\[ a_{n+2} = (p-2)a_{n+1} + (p-1)a_n \]

\end{document}
